\chapter{Conclusion and Future Work}\label{ch:conclusion}

In this thesis an open-source implementation called `Vinculum' of the Lock-Free
Deduplication algorithm described by Gilbert et al.~\cite{Duplicacy} is presented.
It is structured as a reusable library to aid users in applying the algorithm
under real-world conditions such as creating fully functional backup solutions.

While a simple program based on it performed favorably in benchmarks comparing
it to existing backup solutions there is still a lot to do to create a real
alternative to them.
One way users of Vinculum could be supported in this regard is by providing
an additional library which extends the interfaces described in \Cref{sec:traits}
with methods operating on a higher abstraction level than the chunk-based operations
used by Vinculum.
This way it can impose additional constraints and interfaces while preventing all
users of Vinculum from being affected by them.

Other avenues for future work are the investigation of the robustness of the
Lock-Free Deduplication algorithm to the repository experiencing catastrophic
errors such as a power failure or its interaction with different consistency
models besides the one provided by POSIX file systems.
While major providers such as Amazon S3~\cite{s3} and Google Cloud Storage~\cite{GoogleCS}
claim to provide strong consistency investigating the interaction of other consistency
models with the Lock-Free Deduplication algorithm may provide opportunities for
performance improvements.

The source code of Vinculum is available on GitHub\footnote{\url{https://github.com/Deric-W/vinculum}}.
